% s2_background.tex
% Status: DRAFT — generated by TAR-Author; review before locking
% Lock condition: after related work citations verified

\section{Background and Related Work}
\label{sec:background}

\subsection{Catastrophic Forgetting}

Catastrophic forgetting~\cite{McCloskey1989} is the tendency of neural networks to
lose performance on previously learned tasks when trained sequentially on new ones.
The problem arises because stochastic gradient descent updates weights to minimise loss
on the current task without regard for previously acquired knowledge.
In the \emph{task-incremental} setting~\cite{vandeVen2019}, task identity is
provided at test time, which bounds the difficulty by allowing per-task output heads;
shared representations still suffer interference.

\subsection{Regularisation-Based Continual Learning}

Elastic Weight Consolidation (EWC)~\cite{Kirkpatrick2017} penalises changes to
weights identified as important for previous tasks using the diagonal Fisher
Information Matrix as an importance proxy.
The penalty strength $\lambda$ governs the plasticity-stability trade-off: insufficient
$\lambda$ allows forgetting; excessive $\lambda$ prevents learning and can collapse the
network to degenerate predictions.
Our pre-registered EWC $\lambda$-sweep (Appendix~\ref{app:ewc_sweep}) confirms this
sensitivity: $\lambda = 10000$ produces chance-level accuracy on all five test seeds.

Synaptic Intelligence (SI)~\cite{Zenke2017} estimates weight importance online
by accumulating the product of parameter gradients and parameter changes during
training, then applies a quadratic penalty to drift from previous task solutions.
At its published default hyperparameters ($c = 0.1$, $\xi = 0.001$), SI produces
$0.500$ accuracy — chance level for binary classification — on every seed of our
benchmark (Section~\ref{sec:si_collapse}).

\subsection{Replay-Based Methods}

Generative replay~\cite{Shin2017} trains a generative model alongside the
classifier and rehearses synthetic samples of previous tasks.
Example replay~\cite{Rebuffi2017} stores a small episodic memory of real past
samples for rehearsal.
Both families are complementary to the regularisation approach evaluated here; we
do not compare against replay methods because our setting prohibits memory buffers.

\subsection{Thermodynamic Perspectives on Neural Network Training}

Thermodynamic analogies in deep learning have a long history, from energy-based
models to Hopfield networks.
More recently, activation statistics have been studied as proxies for training
dynamics --- large activation standard deviations indicate an exploratory, disordered
regime; small standard deviations indicate convergence.
TCL is, to our knowledge, the first method to use the \emph{ratio} of current to
anchored activation entropy as a \emph{control signal} for continual learning: rather
than measuring importance retrospectively (as EWC and SI do), TCL detects the
consolidation moment as it occurs and triggers weight anchoring at that instant.
